\label{cap:progettazione}

Questo capitolo descrive la progettazione della gestione del ciclo di vita di un Task, a partire dai requisiti definiti nel Capitolo~\ref{cap:requisiti}, ai quali si fa riferimento tramite i relativi identificativi (RF1, RF2, \ldots). La Sezione~\ref{sec:execute} descrive la progettazione di \texttt{execute}, mostrando come un'unica implementazione possa servire sia le invocazioni sincrone sia quelle in streaming. La Sezione~\ref{sec:cancel} tratta la progettazione di \texttt{cancel}, la divisione dei compiti con il framework e il modo in cui \texttt{cancel} ed \texttt{execute} collaborano. La Sezione~\ref{sec:limiti} discute infine i limiti della soluzione e i possibili miglioramenti emersi dall'analisi dell'SDK.



\section{Progettazione di \texorpdfstring{\texttt{execute}}{execute}: modalità stream e invoke}
\label{sec:execute}

Come stabilito nel Capitolo~\ref{cap:requisiti}, l'interfaccia \texttt{AgentExecutor} prevede un unico metodo \texttt{execute}, invocato dal framework indipendentemente dalla modalità scelta dal client. I requisiti RF1 e RF2 richiedono però di gestire correttamente sia le chiamate sincrone (\texttt{SendMessage}) sia quelle in streaming (\texttt{SendStreamingMessage}), e in quest'ultimo caso di restituire anche tutti gli eventi intermedi. Il problema progettuale consisteva quindi nel soddisfare entrambe le modalità con un'unica implementazione di \texttt{execute}.

Una prima soluzione possibile consiste nel far verificare a \texttt{execute}, a ogni invocazione, quale operazione abbia generato la richiesta, scegliendo di conseguenza quale metodo del grafo LangGraph chiamare. Con \texttt{SendMessage}, \texttt{execute} utilizzerebbe \texttt{ainvoke}, un metodo che esegue tutta l'elaborazione e restituisce un unico valore, cioè lo stato finale, solo alla fine. Con \texttt{SendStreamingMessage} utilizzerebbe invece \texttt{astream}, che restituisce un aggiornamento dopo ogni passo del grafo, così da pubblicare i relativi eventi sulla \texttt{EventQueue} mentre l'esecuzione avanza. Si avrebbero quindi due percorsi distinti nella stessa funzione, scelti con un controllo condizionale all'inizio di \texttt{execute}. 

Questo approccio presenta però due limiti concreti. Il primo è che introduce nella logica applicativa dell'agente un'informazione che appartiene al livello di trasporto, cioè quale operazione RPC abbia originato la richiesta, in contrasto con il principio di separazione delle responsabilità. Il secondo è che costringe a mantenere due percorsi di codice paralleli e in gran parte ridondanti, destinati a evolvere in modo indipendente. Ogni caso particolare, come la richiesta di ulteriori input all'utente, andrebbe gestito allo stesso modo in entrambi i percorsi, e un aggiornamento applicato a uno solo dei due renderebbe diverso il comportamento delle due modalità senza che la differenza risulti evidente.

\begin{figure}[htbp]
    \centering
    \begin{tikzpicture}[
        actor/.style={draw, rounded corners, minimum width=2.6cm, minimum height=0.7cm, align=center, fill=blue!5, font=\footnotesize},
        lifeline/.style={dashed, thin},
        msg/.style={-{Stealth[length=2.5mm]}, thick, font=\scriptsize}
    ]
        \node[actor] (c1) at (0,0) {Client};
        \node[actor] (h1) at (4.5,0) {DefaultRequestHandler};
        \node[actor] (e1) at (9,0) {AgentExecutor};
        \draw[lifeline] (c1) -- ++(0,-5.0);
        \draw[lifeline] (h1) -- ++(0,-5.0);
        \draw[lifeline] (e1) -- ++(0,-5.0);

        \draw[msg] (0,-0.8) -- (4.5,-0.8) node[midway, above]{SendMessage};
        \draw[msg] (4.5,-1.5) -- (9,-1.5) node[midway, above]{execute()};
        \draw[dashed, rounded corners] (4.8,-1.9) rectangle (8.7,-3.2);
        \node[font=\tiny, align=center] at (6.75,-2.55) {eventi\\ SUBMITTED / WORKING\\pubblicati ma non inoltrati};
        \draw[msg] (9,-4.0) -- (4.5,-4.0) node[midway, above]{evento COMPLETED};
        \draw[msg] (4.5,-4.6) -- (0,-4.6) node[midway, above]{risposta finale (Task)};

        \node[font=\scriptsize] at (4.5,-5.5) {(a) \texttt{SendMessage} -- il client riceve solo il risultato finale};
    \end{tikzpicture}

    \vspace{1cm}

    \begin{tikzpicture}[
        actor/.style={draw, rounded corners, minimum width=2.6cm, minimum height=0.7cm, align=center, fill=blue!5, font=\footnotesize},
        lifeline/.style={dashed, thin},
        msg/.style={-{Stealth[length=2.5mm]}, thick, font=\scriptsize}
    ]
        \node[actor] (c2) at (0,0) {Client};
        \node[actor] (h2) at (4.5,0) {DefaultRequestHandler};
        \node[actor] (e2) at (9,0) {AgentExecutor};
        \draw[lifeline] (c2) -- ++(0,-6.5);
        \draw[lifeline] (h2) -- ++(0,-6.5);
        \draw[lifeline] (e2) -- ++(0,-6.5);

        \draw[msg] (0,-0.8) -- (4.5,-0.8) node[midway, above]{SendStreamingMessage};
        \draw[msg] (4.5,-1.5) -- (9,-1.5) node[midway, above]{execute()};

        \draw[msg] (9,-2.5) -- (4.5,-2.5) node[midway, above]{evento SUBMITTED};
        \draw[msg] (4.5,-3) -- (0,-3) node[midway, above]{evento SUBMITTED};

        \draw[msg] (9,-4) -- (4.5,-4) node[midway, above]{evento WORKING};
        \draw[msg] (4.5,-4.5) -- (0,-4.5) node[midway, above]{evento WORKING};

        \node[font=\large, fill=white, inner sep=1pt] at (4.5,-5.0) {$\vdots$};

        \draw[msg] (9,-5.5) -- (4.5,-5.5) node[midway, above]{evento COMPLETED};
        \draw[msg] (4.5,-6) -- (0,-6) node[midway, above]{evento COMPLETED};

        \node[font=\scriptsize] at (4.5,-7) {(b) \texttt{SendStreamingMessage} -- ogni evento è inoltrato subito};
    \end{tikzpicture}
    \caption{La stessa esecuzione di \texttt{execute} con due modalità di consegna al client, gestite interamente dal \texttt{DefaultRequestHandler}.}
    \label{fig:sequence-comparison}
\end{figure}

Si è quindi scelto di adottare un'unica strategia di esecuzione, indipendente dalla modalità di invocazione. \texttt{execute} itera sul grafo LangGraph sempre in streaming, tramite \texttt{astream}, e pubblica un evento sulla \texttt{EventQueue} a ogni passo significativo. È poi il framework, e in particolare il \texttt{DefaultRequestHandler}, a stabilire come consumare questa sequenza di eventi \cite{a2apythontutorial} in base alla modalità con cui il client ha effettuato la richiesta, come mostrato in Figura~\ref{fig:sequence-comparison}.

Nel pannello (a) il client invia una richiesta \texttt{SendMessage}. Il gestore delle richieste (\texttt{Default\allowbreak Request\allowbreak Handler}) la riceve e invoca \texttt{execute}, che durante l'esecuzione pubblica sulla \texttt{EventQueue} gli eventi \texttt{SUBMITTED} e \texttt{WORKING}. Il gestore riceve questi eventi ma non li inoltra al client. La connessione resta aperta, senza che vi transiti alcun dato, finché non viene pubblicato un evento che porta il Task in uno stato terminale (\texttt{COMPLETED}, \texttt{FAILED}, \texttt{CANCELED} o \texttt{REJECTED}) oppure interrotto (\texttt{INPUT\_REQUIRED} o \texttt{AUTH\_REQUIRED}) \cite{a2aspec}. Solo a quel punto il gestore restituisce al client un'unica risposta, che contiene il Task. Resta valida, come visto nella Sezione~\ref{sec:a2a-architettura}, la possibilità per il client di chiedere una risposta immediata con \texttt{return\_immediately}.

Nel pannello (b) la stessa esecuzione di \texttt{execute} parte da una richiesta \texttt{SendStreamingMessage}. Anche qui la connessione resta aperta per tutta l'esecuzione, ma il gestore si comporta diversamente. La risposta è uno stream di eventi, e ogni evento pubblicato sulla \texttt{EventQueue} viene inoltrato al client non appena ricevuto. Lo stream si chiude dopo l'evento che porta il Task in uno stato terminale o interrotto.

Durante l'esecuzione il Task passa attraverso una serie di stati, che \texttt{execute} comunica esplicitamente tramite il \texttt{TaskUpdater} e che sono riassunti in Figura~\ref{fig:task-lifecycle} (per brevità, nel diagramma è omesso il prefisso \texttt{TASK\_STATE\_} comune a tutti gli stati).

\begin{figure}[htbp]
    \centering
    \resizebox{\linewidth}{!}{%
    \begin{tikzpicture}[
        state/.style={draw, rounded corners, minimum width=2.8cm, minimum height=1cm, align=center, fill=blue!5, font=\small},
        terminal/.style={state, fill=gray!15, very thick},
        arrow/.style={-{Stealth[length=3mm]}, thick},
        dashedarrow/.style={arrow, dashed},
        dottedarrow/.style={arrow, dotted},
        lbl/.style={font=\footnotesize}
    ]
        \node[state] (submitted) at (-2.5,0) {SUBMITTED};
        \node[state] (working) at (3.5,0) {WORKING};
        \node[state] (input) at (3.5,4) {INPUT\_REQUIRED};
        \node[terminal] (completed) at (10,2) {COMPLETED};
        \node[terminal] (failed) at (10,0) {FAILED};
        \node[terminal] (canceled) at (10,-2) {CANCELED};

        \draw[arrow] (submitted) -- (working) node[midway, above, lbl]{avvio elaborazione};

        \draw[arrow] (3.9,0.5) -- (3.9,3.5) node[midway, right, lbl, xshift=3pt] {input necessario};
        \draw[dashedarrow] (3.1,3.5) -- (3.1,0.5) node[midway, left, lbl, xshift=-3pt] {nuovo messaggio};

        \draw[arrow] (working) to[loop below, looseness=8] node[below=2pt, lbl] {step intermedio} (working);

        \draw[arrow] (working) -- (completed) node[midway, above, sloped, lbl]{successo};
        \draw[arrow] (working) -- (failed) node[midway, above, lbl]{eccezione};
        \draw[arrow] (working) -- (canceled) node[pos=0.4, below, sloped, lbl]{cancellazione};

        \draw[arrow] (submitted.south) -- (-2.5,-2) -- (canceled);
        \node[lbl] at (3.5,-2.3) {cancellazione immediata};

        \draw[dottedarrow] (input.east) -- (12.8,4) -- (12.8,-2) -- (canceled.east);
        \node[lbl, rotate=-90] at (13.15,1) {cancellazione (\texttt{cancel})};
    \end{tikzpicture}%
    }
    \caption{Ciclo di vita del Task. Le frecce continue e tratteggiate rappresentano le transizioni prodotte da \texttt{execute}, quella punteggiata la transizione prodotta direttamente da \texttt{cancel} (Sezione~\ref{sec:cancel}). Gli stati con bordo spesso (in grigio) sono terminali. La specifica A2A prevede anche gli stati \texttt{REJECTED} (terminale) e \texttt{AUTH\_REQUIRED} (interrotto), non utilizzati in questa implementazione \cite{a2aspec}.}
\label{fig:task-lifecycle}
\end{figure}

Quando la richiesta avvia un nuovo Task, \texttt{execute} pubblica per prima cosa un evento che lo porta in \texttt{SUBMITTED}. Subito dopo, prima ancora di iterare sul grafo, pubblica esplicitamente il passaggio a \texttt{WORKING}. In questo modo il Task passa sempre da questo stato prima di arrivare a un esito, anche quando il grafo produce il risultato in un solo passo. Il Task resta in \texttt{WORKING} per tutta l'elaborazione, e a ogni passo intermedio \texttt{execute} pubblica un evento di avanzamento, che in figura corrisponde alla transizione ad anello.

Se l'agente necessita di ulteriori informazioni, il Task passa in \texttt{INPUT\_REQUIRED} e \texttt{execute} termina. Il messaggio successivo dell'utente sullo stesso Task (freccia tratteggiata) porta a una nuova invocazione di \texttt{execute}, che riporta il Task in \texttt{WORKING} e riprende l'elaborazione. In questo caso l'evento \texttt{SUBMITTED} non viene pubblicato di nuovo, perché il Task esiste già. Altrimenti l'esecuzione termina in uno dei tre stati terminali: \texttt{COMPLETED} se va a buon fine, \texttt{FAILED} se si verifica un'eccezione, oppure \texttt{CANCELED} se arriva una richiesta di cancellazione, durante l'elaborazione o, come indica il percorso ad angolo retto, prima ancora che inizi. Il passaggio da \texttt{INPUT\_REQUIRED} a \texttt{CANCELED} (freccia punteggiata) non è prodotto da \texttt{execute}, perché in quello stato non c'è nessuna esecuzione attiva. Se ne occupa direttamente \texttt{cancel}, come descritto nella Sezione~\ref{sec:cancel}.

L'Algoritmo~\ref{alg:execute} formalizza la logica appena descritta. Un aspetto richiede particolare attenzione: la posizione del blocco \texttt{try}, che si apre subito dopo la registrazione del task nel registro dei task in esecuzione. In questo modo qualsiasi cancellazione o eccezione viene intercettata, anche se arriva mentre si stanno pubblicando \texttt{SUBMITTED} o \texttt{WORKING}, e il blocco \texttt{finally} rimuove comunque il task dal registro. Il motivo per cui servono il registro e la gestione di \texttt{CancelledError} è spiegato nella Sezione~\ref{sec:cancel}.

\begin{algorithm}[htbp]
\caption{Logica di \texttt{execute}}
\label{alg:execute}
\begin{algorithmic}[1]
\Function{Execute}{$context$, $event\_queue$}
    \State $query \gets$ \Call{OttieniInputUtente}{$context$}
    \State $task\_id \gets context.task\_id$
    \State $context\_id \gets context.context\_id$
    \State $updater \gets$ nuovo \Call{TaskUpdater}{$event\_queue$, $task\_id$, $context\_id$}
    \State $registro[task\_id] \gets$ task \texttt{asyncio} corrente
    \Try
        \If{$context.current\_task$ non è definito} \Comment{nuovo Task}
            \State Pubblica evento Task($task\_id$, stato = \textsc{Submitted})
        \EndIf
        \State $updater.$\Call{AvviaLavoro}{}
        \For{$step$ in $agente.$\Call{Stream}{$query$, $context\_id$}}
            \If{$step$ indica task completato}
                \State $updater.$\Call{AggiungiArtifact}{$step.risultato$}
                \State $updater.$\Call{Completa}{}
                \State \textbf{break}
            \ElsIf{$step$ richiede input utente}
                \State $updater.$\Call{RichiediInput}{$step.messaggio$}
                \State \textbf{break}
            \Else
                \State $updater.$\Call{AggiornaAvanzamento}{$step.messaggio$}
            \EndIf
        \EndFor
    \Catch{\textsc{CancelledError}}
        \State $updater.$\Call{Cancella}{}
        \State rilancia \textsc{CancelledError}
    \Catch{\textsc{Exception} $e$}
        \State $updater.$\Call{Fallito}{}
        \State solleva \textsc{InternalError}
    \Finally
        \State rimuovi $task\_id$ dal registro
    \EndTry
\EndFunction
\end{algorithmic}
\end{algorithm}

\FloatBarrier


\section{Progettazione di \texorpdfstring{\texttt{cancel}}{cancel}}
\label{sec:cancel}

Come stabilito nel Capitolo~\ref{cap:requisiti}, l'interfaccia \texttt{AgentExecutor} definisce \texttt{cancel} come un metodo asincrono basato su \texttt{asyncio}. In questo contesto, invocare \texttt{cancel()} su un task \texttt{asyncio} non ne interrompe subito l'esecuzione, ma programma il sollevamento di un'eccezione \texttt{asyncio.CancelledError} al primo punto di sospensione utile, tipicamente un'istruzione \texttt{await} \cite{pythonasynciotask}. Si parla quindi di cancellazione \emph{cooperativa}, perché è il codice in esecuzione, e non il chiamante, a determinare quando l'interruzione avviene davvero.

\subsection{Divisione dei compiti con il framework}

Prima di progettare \texttt{cancel} è necessario chiarire quali operazioni siano già svolte dall'SDK. Nella versione adottata, il \texttt{DefaultRequestHandler} associa a ogni Task un task \texttt{asyncio} dedicato, che nel seguito verrà indicato come \emph{task di esecuzione}, e che invoca \texttt{execute} a ogni nuova richiesta. Il task di esecuzione non termina insieme a \texttt{execute}, ma resta attivo tra una richiesta e l'altra in attesa di nuovi messaggi, e si chiude solo quando il Task arriva in uno stato terminale. Di conseguenza, esso risulta ancora attivo anche quando il Task si trova in \texttt{INPUT\_REQUIRED}.

Quando arriva una richiesta \texttt{CancelTask}, il gestore esegue quattro operazioni in quest'ordine. Innanzitutto verifica che l'identificativo corrisponda a un Task esistente, e in caso contrario solleva \texttt{TaskNotFoundError}. Poi, se il task di esecuzione è ancora attivo, invoca il metodo \texttt{cancel} dell'\texttt{AgentExecutor} e, solo quando questo è terminato, cancella il task di esecuzione. Infine, se il Task non è ancora in uno stato terminale, imposta il suo stato memorizzato a \texttt{CANCELED} \cite{a2apythonsdk}.

L'ordine di queste operazioni non è casuale, ed è motivato esplicitamente nel codice dell'SDK. Il task di esecuzione viene cancellato solo dopo che il metodo \texttt{cancel} dell'\texttt{AgentExecutor} è terminato. Così il componente responsabile dello stato terminale del Task può ancora pubblicarlo sulla coda di eventi, che altrimenti rischierebbe di essere chiusa da una cancellazione anticipata del task. Lo stato \texttt{CANCELED} scritto direttamente dal framework, invece, non ha nessuna garanzia di arrivare a un client in streaming, perché quando viene registrato le code di eventi potrebbero essere già chiuse \cite{a2apythonsdk}. L'SDK presuppone pertanto che sia l'agente a gestire lo stato terminale, in linea con il contratto dell'interfaccia, secondo cui l'agente deve provare a interrompere il Task e pubblicare l'evento che lo porta in \texttt{CANCELED} \cite{a2apythondocs}. L'intervento automatico del framework costituisce quindi un meccanismo di \emph{fallback}, e non un'alternativa alla corretta gestione da parte dell'agente.

Dall'analisi del codice dell'SDK emerge inoltre che, nella versione adottata, il comportamento verso i Task già terminati non coincide in tutti i casi con quello previsto dalla specifica. Questi controlli però avvengono nel gestore delle richieste, prima che la richiesta arrivi all'\texttt{AgentExecutor}, e non possono essere corretti modificando \texttt{execute} e \texttt{cancel}. La questione è quindi ripresa tra i limiti della soluzione, nella Sezione~\ref{sec:limiti}.

\subsection{Il registro dei task in esecuzione}

Resta da definire il compito di \texttt{cancel} nell'\texttt{AgentExecutor}. A tal fine è stato introdotto un \emph{registro dei task in esecuzione}, cioè un dizionario che associa a ogni \texttt{task\_id} il task \texttt{asyncio} che sta eseguendo \texttt{execute}. Il registro viene popolato da \texttt{execute} all'avvio e ripulito dalla stessa funzione al termine, qualunque sia l'esito (Algoritmo~\ref{alg:execute}). La presenza di un Task nel registro indica quindi l'esistenza di un'esecuzione attiva di \texttt{execute}. Dato che il registro è indicizzato per \texttt{task\_id}, le esecuzioni di Task diversi restano del tutto indipendenti (requisito RF7).

Nella versione adottata il registro non è strettamente necessario per interrompere l'esecuzione. Il task \texttt{asyncio} in cui viene eseguito \texttt{execute} è infatti lo stesso task di esecuzione gestito dall'SDK, che, come descritto in precedenza, viene comunque cancellato dopo l'invocazione di \texttt{cancel} \cite{a2apythonsdk}. La chiamata a \texttt{cancel()} sul task registrato è quindi ridondante. Si è scelto comunque di mantenerla, perché il contratto dell'interfaccia assegna esplicitamente all'agente il compito di provare a interrompere il Task \cite{a2apythondocs}, mentre la cancellazione forzata da parte del gestore delle richieste è un dettaglio implementativo che l'interfaccia non garantisce e che potrebbe cambiare nelle versioni successive dell'SDK. La ridondanza, peraltro, non produce effetti collaterali, perché più richieste di cancellazione rivolte allo stesso task \texttt{asyncio} prima che questo riprenda l'esecuzione portano al sollevamento di un'unica \texttt{CancelledError} \cite{pythonasynciotask}.

La funzione principale del registro è però un'altra: permettere a \texttt{cancel} di stabilire se debba pubblicare direttamente l'evento \texttt{CANCELED}. Il gestore delle richieste invoca \texttt{cancel} sia quando \texttt{execute} è in corso sia quando il Task è in attesa di input, perché in entrambi i casi il suo task di esecuzione risulta ancora attivo. Le due situazioni vanno però gestite in modo opposto. Se \texttt{execute} è in corso, \texttt{cancel} non deve pubblicare \texttt{CANCELED}, perché se ne occupa già \texttt{execute}, e una doppia pubblicazione genererebbe eventi duplicati sullo stream. Se invece il Task è in attesa di input, non c'è nessuna esecuzione che possa pubblicare l'evento, quindi deve farlo \texttt{cancel}. Altrimenti il framework scriverebbe \texttt{CANCELED} solo nello stato memorizzato, e un client iscritto agli aggiornamenti del Task tramite \texttt{SubscribeToTask} vedrebbe lo stream chiudersi senza ricevere alcun evento terminale; si tratta di un limite dell'SDK già segnalato nel suo repository \cite{a2apythonissue1175}. Soltanto \texttt{execute} sa se un'esecuzione è in corso, e il registro serve a rendere questa informazione disponibile a \texttt{cancel}.

In quest'ultimo caso conta anche il tipo di evento pubblicato. \texttt{cancel} deve usare il \texttt{TaskUpdater}, che produce un \texttt{TaskStatusUpdateEvent}, e non pubblicare un nuovo oggetto Task con stato \texttt{CANCELED}: se riceve un Task con un identificativo già esistente, l'SDK lo ignora e non ne aggiorna lo stato \cite{a2apythonsdk}. Per lo stesso motivo \texttt{execute} pubblica il Task iniziale solo quando il Task è nuovo, e non quando riprende dopo una richiesta di input (Algoritmo~\ref{alg:execute}).

Combinando i controlli dell'SDK e il registro, una richiesta di cancellazione può ricadere in tre scenari, ciascuno legato a uno o più requisiti del Capitolo~\ref{cap:requisiti}:

\begin{enumerate}
    \item \textbf{Task in esecuzione} (RF3). Il Task è presente nel registro, quindi \texttt{cancel} invoca \texttt{cancel()} sul task \texttt{asyncio} corrispondente, come ulteriore garanzia oltre all'interruzione che l'SDK esegue comunque, e non pubblica nessun evento. Lo stato \texttt{CANCELED} viene pubblicato da \texttt{execute}, come discusso più avanti.
    \item \textbf{Task in attesa, senza esecuzione attiva} (RF4). Il Task esiste, non è in uno stato terminale e non è nel registro. È il caso tipico di un Task in \texttt{INPUT\_REQUIRED}, per il quale \texttt{execute} è già terminata in attesa di un nuovo messaggio (Sezione~\ref{sec:execute}). Non essendoci nulla da interrompere, \texttt{cancel} pubblica direttamente l'evento che porta il Task in \texttt{CANCELED}.
    \item \textbf{Task non cancellabile} (RF5, RF6). Se il Task non esiste, il gestore delle richieste dell'SDK rifiuta la richiesta prima che arrivi all'\texttt{AgentExecutor}, sollevando \texttt{TaskNotFoundError} (RF5). Lo stesso avviene, di norma, se il Task è già in uno stato terminale, con \texttt{TaskNotCancelableError}. In entrambi i casi la richiesta non raggiunge \texttt{cancel}, per cui le richieste di cancellazione successive alla prima lasciano il Task in \texttt{CANCELED} (RF6). L'eccezione a questo comportamento è discussa nella Sezione~\ref{sec:limiti}.
\end{enumerate}

La Figura~\ref{fig:cancel-flow} mostra l'ordine in cui vengono valutate queste condizioni e quale componente si occupa di ciascun controllo, mentre l'Algoritmo~\ref{alg:cancel} formalizza la logica di \texttt{cancel}.

\begin{figure}[htbp]
    \centering
    \resizebox{\linewidth}{!}{%
    \begin{tikzpicture}[
        font=\normalsize,
        base/.style={draw, rounded corners, align=center, minimum height=1cm, inner sep=6pt},
        process/.style={base, fill=blue!5, minimum width=3cm},
        decision/.style={base, fill=yellow!10, minimum width=3cm},
        stop/.style={base, fill=gray!15, very thick, minimum width=3cm},
        arrow/.style={-{Stealth[length=3mm]}, thick},
        lbl/.style={font=\small, fill=white, inner sep=1.5pt},
        group/.style={draw, dashed, rounded corners, gray},
        grouplbl/.style={font=\small\itshape, anchor=north west, gray!70!black}
    ]
        \draw[group] (-8.3,-1.3) rectangle (2.3,-6.0);
        \node[grouplbl] at (-8.2,-1.35) {SDK (\texttt{DefaultRequestHandler})};
        \draw[group] (-6.6,-6.5) rectangle (6.6,-11.0);
        \node[grouplbl] at (-6.5,-6.55) {\texttt{BaseExecutor.cancel}};
        
        \node[process]  (start) at (0,0)     {\textsc{CancelTask}($task\_id$)};
        \node[decision] (q1)    at (0,-2.6)  {Task\\memorizzato?};
        \node[stop]     (err1)  at (-5.4,-2.6) {Solleva\\\textsc{TaskNotFoundError}};
        \node[decision] (q2)    at (0,-4.8)  {Stato\\terminale?};
        \node[stop]     (err2)  at (-5.4,-4.8) {Solleva\\\textsc{TaskNotCancelableError}};
        \node[decision] (q3)    at (0,-7.6)  {$task\_id$ nel\\registro?};
        \node[process]  (upd)   at (-3.8,-9.9) {Pubblica evento\\\textsc{Canceled}};
        \node[process]  (canc)  at (3.8,-9.9)  {$task\_asyncio$.\textsc{Cancel}()};
        \node[stop]     (end1)  at (-3.8,-12.3) {Fine\\(scenario 2)};
        \node[stop]     (end2)  at (3.8,-12.3)  {Fine (scenario 1):\\\textsc{Canceled} pubblicato\\da \texttt{execute}};
        
        \draw[arrow] (start) -- (q1);
        \draw[arrow] (q1) -- node[lbl]{no} (err1);
        \draw[arrow] (q1) -- node[lbl]{sì} (q2);
        \draw[arrow] (q2) -- node[lbl]{sì} (err2);
        \draw[arrow] (q2) -- node[lbl, pos=0.3]{no} (q3);
        \draw[arrow] (q3) -- node[lbl, pos=0.45]{no} (upd);
        \draw[arrow] (q3) -- node[lbl, pos=0.45]{sì} (canc);
        \draw[arrow] (upd) -- (end1);
        \draw[arrow] (canc) -- (end2);
    \end{tikzpicture}%
    }
    \caption{Gestione di una richiesta \texttt{CancelTask}. I controlli sull'esistenza e sullo stato del Task vengono eseguiti dall'SDK prima di invocare \texttt{BaseExecutor.cancel}, che distingue i due scenari in base al registro dei task in esecuzione. Dopo \texttt{BaseExecutor.cancel}, l'SDK cancella a sua volta il task di esecuzione, se è ancora attivo.}
    \label{fig:cancel-flow}
\end{figure}

\begin{algorithm}[htbp]
\caption{Logica di \texttt{cancel}}
\label{alg:cancel}
\begin{algorithmic}[1]
\Function{Cancel}{$context$, $event\_queue$} \Comment{\texttt{BaseExecutor}}
    \State $task\_id \gets context.task\_id$
    \If{$task\_id$ è presente nel registro} \Comment{scenario 1}
        \State $registro[task\_id].$\Call{Cancel}{}
        \State \Return
    \EndIf
    \State $updater \gets$ nuovo \Call{TaskUpdater}{$event\_queue$, $task\_id$, $context.context\_id$}
    \State $updater.$\Call{Cancella}{} \Comment{scenario 2}
\EndFunction
\end{algorithmic}
\end{algorithm}


\subsection{Responsabilità dello stato terminale}

Per lo scenario~1 si è adottato un principio preciso, secondo cui il ciclo di vita di un Task debba essere gestito da chi ne sta eseguendo l'elaborazione. Per questo \texttt{cancel} non porta mai autonomamente in \texttt{CANCELED} un Task in esecuzione. Si limita a richiedere la cancellazione del task \texttt{asyncio} corrispondente, richiesta che, nella versione adottata, l'SDK ripete subito dopo. È \texttt{execute}, nel blocco \texttt{except} dedicato a \texttt{CancelledError} (Algoritmo~\ref{alg:execute}), a intercettare l'eccezione, pubblicare l'evento che porta il Task in \texttt{CANCELED} ed eseguire le operazioni di pulizia necessarie. Le responsabilità sono state separate in questo modo perché solo \texttt{execute}, nel punto esatto in cui l'elaborazione viene interrotta, conosce ciò che è stato effettivamente completato fino a quel momento. La scelta è inoltre coerente con l'ordine delle operazioni dell'SDK descritto in precedenza. Dato che \texttt{execute} pubblica \texttt{CANCELED} prima di terminare, e quindi prima che le code di eventi vengano chiuse, anche un client in streaming riceve l'evento terminale.

Per lo stesso motivo, dopo aver aggiornato lo stato del Task, il blocco \texttt{except} di \texttt{execute} deve rilanciare \texttt{CancelledError}. Se l'eccezione venisse soppressa, il task \texttt{asyncio} risulterebbe terminato normalmente e non cancellato. Lo stato applicativo del Task sarebbe correttamente \texttt{CANCELED}, mentre per l'event loop il task si sarebbe concluso con successo, e i meccanismi di \texttt{asyncio} che si basano sulla corretta propagazione della cancellazione non funzionerebbero come previsto \cite{pythonasynciotask}.

Il comportamento descritto resta corretto finché i controlli di \texttt{cancel} e le transizioni di \texttt{execute} non si sovrappongono nel tempo. I casi in cui questo non avviene sono discussi nella sezione seguente.

\FloatBarrier

% =====================================================================
\section{Limiti e possibili miglioramenti}
\label{sec:limiti}

L'analisi del codice dell'SDK e del modello di esecuzione di \texttt{asyncio} ha messo in luce alcuni casi che la soluzione descritta non copre. Si tratta di situazioni che dipendono da tempi molto stretti, oppure da componenti esterni all'\texttt{AgentExecutor}, e che non sono state riprodotte nei test descritti nel Capitolo~\ref{cap:validazione}. Per ciascuna si descrivono il problema e una possibile soluzione, che costituisce uno sviluppo futuro del lavoro.

\subsection{Operazioni su Task già terminati}

La specifica prevede che una richiesta \texttt{CancelTask} su un Task già terminato venga respinta con \texttt{TaskNotCancelableError}, e che un nuovo messaggio inviato a un Task terminato venga respinto con \texttt{UnsupportedOperationError} \cite{a2aspec}. Nella versione adottata dell'SDK il primo errore viene sollevato solo se il gestore delle richieste ha già rilasciato le risorse associate al Task: se la richiesta arriva subito dopo il completamento, non produce nessun errore e il Task viene restituito nel suo stato corrente, per cui il risultato dipende dal momento in cui arriva la richiesta. Il secondo caso viene invece respinto con \texttt{InvalidParamsError}, un comportamento che è oggetto di una modifica proposta nel repository dell'SDK dopo la versione adottata \cite{a2apythonsdk,a2apythonpr1268}.

Entrambi i controlli avvengono nel gestore delle richieste, prima che la richiesta arrivi all'\texttt{AgentExecutor}. Una possibile soluzione consiste nell'estendere il \texttt{DefaultRequestHandler} con una sottoclasse che, prima di passare la richiesta all'implementazione dell'SDK, legga lo stato memorizzato del Task e, se è terminale, sollevi l'errore previsto dalla specifica. Il controllo e la delega non avverrebbero comunque in modo atomico: se il Task diventasse terminale proprio tra le due operazioni, prevarrebbe il comportamento dell'SDK.

\subsection{Cancellazione durante la pubblicazione dell'esito}

In \texttt{asyncio} l'esecuzione è concorrente ma non parallela: in ogni istante è attiva al massimo una coroutine, e il controllo passa da una all'altra solo nei punti di sospensione, cioè quando la coroutine in esecuzione resta in attesa tramite \texttt{await} \cite{pythonasynciotask}. Ogni \texttt{await} che sospende davvero \texttt{execute} è quindi un punto in cui una cancellazione può intervenire, sollevando \texttt{CancelledError}.

Questo vale anche per la pubblicazione dell'esito. Quando il grafo ha prodotto il risultato, \texttt{execute} lo pubblica con \textsc{AggiungiArtifact} e \textsc{Completa}, oppure con \textsc{RichiediInput}. Queste operazioni non sono istantanee: il \texttt{TaskUpdater} serializza gli aggiornamenti di stato con un lock, e pubblicare un evento significa inserirlo in una coda \texttt{asyncio} di capacità limitata \cite{a2apythonsdk}. Se una richiesta di cancellazione arriva in questa breve finestra, \texttt{CancelledError} viene sollevata durante la pubblicazione, il blocco \texttt{except} pubblica \texttt{CANCELED} e l'esito già prodotto dal grafo può andare perso. Il Task risulterebbe così cancellato anche se l'elaborazione era di fatto conclusa.

Una possibile soluzione combina due accorgimenti. Il primo è una variabile, impostata appena il grafo produce l'esito e prima di pubblicarlo, che impedisca al blocco \texttt{except} di sovrascriverlo con \texttt{CANCELED}. Il secondo è proteggere la pubblicazione con \texttt{asyncio.shield}, che esegue un'operazione in un task separato e ne impedisce l'interruzione da parte di una cancellazione. Poiché \texttt{shield} solleva comunque subito \texttt{CancelledError} nel chiamante, l'attesa andrebbe ripetuta fino al termine della pubblicazione, e solo dopo la cancellazione andrebbe propagata \cite{pythonasynciotask}. In questo modo l'esito verrebbe sempre accodato prima che \texttt{execute} termini, e quindi prima che l'SDK chiuda le code di eventi.

\subsection{Cancellazione tra la fine di \texorpdfstring{\texttt{execute}}{execute} e l'aggiornamento dello stato}

Resta infine una finestra che non dipende dalla logica di \texttt{execute}, ma dal fatto che la pubblicazione di un evento e il suo effetto sullo stato memorizzato sono separati, e il secondo è gestito dal framework. Nella versione adottata, l'SDK considera concluso un Task solo dopo averne elaborato l'evento terminale \cite{a2apythonsdk}. Se una richiesta di cancellazione arriva dopo che \texttt{execute} si è rimossa dal registro, ma prima che l'SDK abbia elaborato l'esito appena pubblicato, \texttt{cancel} osserva un Task non terminale e senza esecuzione attiva. La richiesta ricade quindi nello scenario~2, e l'evento \texttt{CANCELED} può essere pubblicato dopo l'esito.

Dato che sia la lettura sia l'aggiornamento dello stato memorizzato avvengono fuori dall'\texttt{AgentExecutor}, questa finestra non si può chiudere modificando solo \texttt{execute} e \texttt{cancel}. Per una soluzione completa servirebbe un aggiornamento condizionato dello stato memorizzato, che rifiuti qualsiasi transizione a partire da uno stato terminale. Tale estensione riguarda lo strato di persistenza del framework e costituisce un possibile sviluppo futuro.
